% !编译配方：xelatex*2
% !请务必选择该配方，若无该配方，请手动编译两次xelatex或自行定义该编译配方

% #region 导言区
\documentclass[a4paper,12pt,fontset=none,UTF8]{ctexbook}
\input{settings/usedPackages.tex}       % #载入宏包
\input{settings/baseSettings.tex}       % #载入文档基础配置
% !设置自定义功能
\newcommand{\vect}[1]{\symbf{#1}} % 粗体向量
\newcommand{\overtext}[2]{%
    \begin{tabular}[b]{@{}c@{}}
        \textsf{\textbf{\fontsize{5pt}{5pt}\selectfont #2}} \\[-1.4ex]
        \textbf{#1}
    \end{tabular}%
}

         % #载入自定义命令
% #endregion


\begin{document}
\makemycover{量子力学}{Quantum Mechanics}
\makecontents
% --------正文--------

\chapter{基本概念}
\section{何谓量子}
\begin{center}
    “场的最小激发是量子。” \\
    “量子代表了一种观点，即：‘世界不是连续的，而是分立的’。”
\end{center}

\section{一点历史}
\subsection{黑体辐射}
使用经典理论给出公式的Wein、Rayleigh和Jeans都失败了。
Rayleigh-Jeans公式以能量按自由度均分定理为基础：
\begin{equation}
    I(\omega) = \dfrac{\omega^2}{\pi^2 c^2}kT
\end{equation}
在$\omega$较大时快速发散，即\textcolor{PurpleDeep}{\textbf{紫外灾难}}。

Planck以插值法通过Wien公式和Rayleigh-Jeans公式拟合得到了Plunck公式。
\begin{derivation}
    \[
        I(w) = \dfrac{1}{\textcolor{PurpleDeep}{e^{a\omega/kT}}-1} \cdot \dfrac{b\omega^3}{\pi^2 c^2}\quad a,b\text{为常数}
    \]
    Planck发现$\textcolor{PurpleDeep}{e^{a\omega/kT}}$形似Bolztmann分布$N_n = N_1 e^{-(E_n-E_1)/kT}$。
    
    因此，他假设\textbf{辐射场中偶极振子的能量是一份一份的，每份能量依赖于角频率}。
    \begin{equation*}
        E = a\omega \qquad N_n = N_0 e^{-E_n/kT} = N_0 e^{-na\omega/kT}
    \end{equation*}
    平均能量：
    \begin{align*}
        \langle E \rangle &= \frac{\sum_{n} N_n \cdot E_n}{\sum_{n} N_n} \\[1ex]
        &= \frac{\sum_{n} N_0 e^{-na\omega/kT} \cdot na\omega}{\sum_{n} N_0 e^{-na\omega/kT}} \\[1ex]
        &= \frac{a\omega (0 \cdot e^{-0/kT} + 1 \cdot e^{-1a\omega/kT} + 2 \cdot e^{-2a\omega/kT} + \cdots)}{e^{-0/kT} + e^{-1a\omega/kT} + e^{-2a\omega/kT} + \cdots} \\[1ex]
        &\xlongequal{\varepsilon = e^{-a\omega/kT}} \frac{a\omega (\varepsilon + 2\varepsilon^2 + 3\varepsilon^3 + \cdots)}{1 + \varepsilon + \varepsilon^2 + \varepsilon^3 + \cdots} \\[1ex]
        &= \frac{a\omega \varepsilon \frac{d}{d\varepsilon} (1 + \varepsilon + \varepsilon^2 + \varepsilon^3 + \cdots)}{1 + \varepsilon + \varepsilon^2 + \varepsilon^3 + \cdots} \\[1ex]
        &\xlongequal{\text{Taylor 展开}} \frac{a\omega \varepsilon \frac{d}{d\varepsilon} (\frac{1}{1-\varepsilon})}{\frac{1}{1-\varepsilon}} \\[1ex]
        &= a\omega \varepsilon \frac{1}{1-\varepsilon} = a\omega \frac{e^{-a\omega/kT}}{1 - e^{-a\omega/kT}} \\[1ex]
        &= a\omega \frac{1}{e^{a\omega/kT} - 1}
    \end{align*}
    Rayleigh-Jeans 公式中的能量项为$kT$，用$\langle E \rangle$取代之：
    \[
    I(\omega) = \frac{\omega^2}{\pi^2 c^2} \textcolor{PurpleDeep}{kT} = \frac{\omega^2}{\pi^2 c^2} \cdot \textcolor{PurpleDeep}{a\omega \frac{1}{e^{a\omega/kT} - 1}} = \frac{1}{e^{a\omega/kT} - 1} \frac{a\omega^3}{\pi^2 c^2}
    \]
    这就是Plunck公式，式中$a$即为如今的$\hbar $。
\end{derivation}

\subsection{光电效应}
Einstein基于Planck的假说，提出光是一份一份的\overtext{能量子}{energy quantum}/\overtext{光量子}{light quantum}构成的，每个光子的能量$E=h\nu=\hbar\omega$。
\[
    E^2=m^2c^4=m_0^2c^4+(pc)^2 \xlongequal{\text{光量子}m_0=0}p^2c^2 \Rightarrow  E = pc \Rightarrow p = \frac{h\mu}{c}=\frac{h}{\lambda}
\]

\subsection{Comptom效应}

\subsection{de Brogile波}

\section{Schrödinger方程}
\subsection{Schrödinger方程的推导}
\begin{derivation}
    Schrödinger根据光的波动性和de Brogile的波粒二象性猜测出了实物粒子的波动方程：
    \[\varphi(x,t) = A\cos(kx-\omega t)\]
    根据Euler公式$e^{i\theta} = \cos\theta + i\sin\theta$有：
    \[\varphi(x,t) = A\operatorname{Re}\{e^{i(kx-\omega t)}\}\]
    为了便于数学处理，用整个复数表示波：
    \[\varphi(x,t) = Ae^{i(kx-\omega t)}\]
    实物粒子$E = \hbar\omega \Rightarrow k = \frac{2\pi}{\lambda} = \frac{p}{\hbar}$代入上式中，于是有自由实物粒子的波动方程：
    \[\psi(x,t) = Ae^{\frac{i}{\hbar}(px-Et)}\]
    考察波函数关于t的变化：
    \[\frac{\partial}{\partial t}\psi = -\frac{i}{\hbar}EAe^{\frac{i}{\hbar}(px-Et)} \Rightarrow \textcolor{PurpleDeep}{i\hbar\frac{\partial}{\partial t}\psi=E\psi}\]
    考察波函数关于x的变化，为了获取能量项，取二阶导：
    \[\frac{\partial^2}{\partial x^2}\psi = \left(-\frac{p}{i\hbar}\right)^2 Ae^{\frac{i}{\hbar}(px-Et)} = \left(-\frac{1}{i\hbar}\right)^2 2\mu E\psi \Rightarrow \textcolor{PurpleDeep}{\frac{1}{2\mu}\left(-i\hbar\frac{\partial}{\partial x}\right)^2 \psi = E\psi}\]
    联立上述二式：
    \[i\hbar\frac{\partial}{\partial t}\psi = \frac{1}{2\mu}\left(-i\hbar\frac{\partial}{\partial x}\right)^2 \psi = E\psi\]
    这就是Schrödinger方程
\end{derivation}
事实上，Schrödinger作上述推导\textbf{纯粹基于猜想}，没有严密的数学推导，它在量子力学中是作为公理存在的。

\subsection{Born统计诠释}
\begin{center}
    \textit{波函数$\psi$的意义是什么？}
    
    \textit{到底是什么在波动？}
\end{center}

Born类比光波的模平方等于光强，即光子在该处的密度大小，故而对实物粒子，\textbf{波函数的模平方表示实物粒子出现在该处的概率密度大小}

因此有\textbf{归一化条件}：
\[\int_{-\infty }^{+\infty}|\psi|^2 \dd{\tau } = 1\]

而波函数本身是概率振幅，是概率的平方根，没有实际意义。

\section{电子双缝干涉实验}
\subsection{电子是如何穿过双缝的？}
电子以概率波的形式同时穿过两缝并产生干涉，落到光屏上又坍缩为粒子。
\begin{derivation}
    设电子通过缝$1$的状态为$\psi_1$，通过缝$2$的状态为$\psi_2$，这种某个量确定的态称为\textbf{本征态}。

    \noindent \textbf{若不施以观测}：

    通过双缝后的粒子处于叠加态:
    \[\psi = c_1\psi_1 + c_2\psi_2\text{(概率振幅叠加)} \qquad \text{归一化条件：}|c_1|^2 + |c_2|^2 = 1\]

\end{derivation}

\end{document}